\documentclass[11pt]{article}

% ---------------------------------------------------------------
% Packages
% ---------------------------------------------------------------
\usepackage[margin=1in]{geometry}
\usepackage{amsmath,amssymb,amsthm}
\usepackage{enumitem}
\usepackage{xcolor}
\usepackage{hyperref}

% ---------------------------------------------------------------
% Macros
% ---------------------------------------------------------------
\newcommand{\cA}{\mathcal{A}}
\newcommand{\cD}{\mathcal{D}}
\newcommand{\cS}{\mathcal{S}}
\newcommand{\cN}{\mathcal{N}}
\newcommand{\R}{\mathbb{R}}
\newcommand{\Z}{\mathbb{Z}}
\newcommand{\E}{\mathbb{E}}
\newcommand{\V}{\mathbb{V}}
\newcommand{\Lap}{\mathrm{Lap}}
\newcommand{\FPR}{\mathrm{FPR}}
\newcommand{\TPR}{\mathrm{TPR}}
\newcommand{\eps}{\varepsilon}
\newcommand{\user}{\mathrm{user}}

% ---------------------------------------------------------------
% Answer block
% ---------------------------------------------------------------
\newenvironment{answer}{\par\medskip\noindent\textbf{Answer.}\ }{\par\bigskip}

% Problem headings
\newcommand{\problem}[2]{\section*{#1 \hfill\normalsize(#2)}}

\setlist[enumerate,1]{label=(\alph*), leftmargin=*, itemsep=6pt}
\setlist[enumerate,2]{label=(\roman*), leftmargin=*, itemsep=4pt}

% ---------------------------------------------------------------
\begin{document}

\begin{center}
  {\LARGE\bfseries AI Privacy: Theory and Practice}\\[6pt]
  {\Large Homework 1: Introduction to Differential Privacy (Lecture 1)}\\[8pt]
  Instructor: Eli Chien\\
  Department of Electrical Engineering, National Taiwan University\\[4pt]
  Fall 2026 \qquad Due: October 12, 2026
\end{center}

\vspace{4pt}
\noindent\text{Name:} Ying Pei Lin \\
\text{Student ID:} R15942158

% ===============================================================
\problem{Problem 1: Properties of the DP definition}{25 points}
% ===============================================================

This problem fills in three gaps that the lecture left as exercises, and then
makes the hypothesis-testing interpretation of DP quantitative.

\begin{enumerate}

\item (4 pts) \textbf{One-sided suffices.} Remark 2 in the lecture says that,
since the roles of $D$ and $D'$ are symmetric under both add/remove and
replacement adjacency, it is enough to require the one-sided bound
$\Pr(\cA(D)\in S)\le e^{\eps}\Pr(\cA(D')\in S)+\delta$. Explain precisely
why. Then give an example of a non-symmetric adjacency relation (one where
$D\simeq D'$ does not imply $D'\simeq D$) and explain what could go wrong if
one only imposed the one-sided bound for it.

\begin{answer}

\textbf{Why one side suffices.}

Under add/remove adjacency, one datasets contains one extra record:

\[
  D \simeq D' \iff D = D' \cup \{x\}
  \ \text{or}\ D = D' \setminus \{x\}
  \ \text{for some } x \in \mathcal{X}
\]

Exchanging the roles of $D$ and $D'$:

\[
\begin{aligned}
D' \simeq D
&\iff D' = D \cup \{x\} \ \text{or}\ D' = D \setminus \{x\}
  \ \text{for some } x \in \mathcal{X} \\
&\iff D = D' \setminus \{x\} \ \text{or}\ D = D' \cup \{x\}
  \ \text{for some } x \in \mathcal{X} \\
&\iff D \simeq D'
\end{aligned}
\]

Under replacement adjacency, the datasets have the same size and differ in one record:

\[
  D \simeq D' \iff D = E\cup\{a\},\ D' = E\cup\{b\}
  \text{ for some } E \text{ and records } a\neq b
\]

Similarly,

\[
\begin{aligned}
D' \simeq D
&\iff D' = E\cup\{a\},\ D = E\cup\{b\}
  \ \text{for some } E \text{ and records } a\neq b \\
&\iff D = E\cup\{b\},\ D' = E\cup\{a\}
  \ \text{for some } E \text{ and records } b\neq a \\
&\iff D \simeq D'
\end{aligned}
\]

Suppose the one-sided bound

\[
  \Pr(\cA(D)\in S)\le e^{\eps}\Pr(\cA(D')\in S)+\delta
\]

holds for all adjacent pairs $(D, D')$.
Since both adjacency relations are symmetric, $(D', D)$ is also adjacent.
Applying the same bound to $(D', D)$ gives

\[
  \Pr(\cA(D')\in S)\le e^{\eps}\Pr(\cA(D)\in S)+\delta .
\]

Therefore the bound holds in both directions, and requiring only one side
is enough.

\textbf{A non-symmetric relation.}

Consider add-only adjacency:

\[
 D \simeq D' \iff D = D' \cup \{x\} \ \text{for some } x \in \mathcal{X}.
\]

Here $D$ is always the larger dataset, so the reverse direction is never checked.
For example, fix a record $z$, let $\eps=\ln 2$ and $\delta=0$, and define

If $z\in D$, then
\[
\Pr(\cA(D)=0)=1,
\qquad
\Pr(\cA(D)=1)=0.
\]

If $z\notin D$, then
\[
\Pr(\cA(D)=0)=\frac{1}{2},
\qquad
\Pr(\cA(D)=1)=\frac{1}{2}.
\]

If $x=z$, then $z\in D$ and $z\notin D'$. Since $e^\eps=2$,

\[
\Pr(\cA(D)=0)
=1
=2\cdot\frac{1}{2}
=e^\eps\Pr(\cA(D')=0),
\qquad
\Pr(\cA(D)=1)=0.
\]

Thus, the one-sided bound holds. However, in the reverse direction, for
$S=\{1\}$,

\[
\Pr(\cA(D')=1)
=\frac{1}{2}
>
0
=e^\eps\Pr(\cA(D)=1).
\]

Therefore, the reverse bound fails. Moreover, observing output $1$ reveals
with certainty that $z$ is not in the dataset. Hence, for a non-symmetric
adjacency relation, imposing only the one-sided bound may fail to protect
information in the reverse direction.

\end{answer}

\item (6 pts) \textbf{Post-processing with a randomized map.} The lecture
proved the post-processing theorem for a deterministic $g:\cS\mapsto\cS'$
and left the randomized case as homework. Let $\cA$ be $(\eps,\delta)$-DP and
let $g$ be a randomized map whose randomness is independent of $\cA$ and of
the data. Prove that $g\circ\cA$ is $(\eps,\delta)$-DP.

\textit{Hint.} Write $g(s)=G(s,R)$ for a deterministic $G$ and an independent
random seed $R$, condition on $R$, and apply the deterministic result.

\begin{answer}

\end{answer}

\item (7 pts) \textbf{Sequential (adaptive) composition.} The lecture proved
simple composition, where $\cA_1$ and $\cA_2$ are run independently on $D$,
and asked what happens in the sequential case where the output of $\cA_1$ is
fed into $\cA_2$. Formalize this as follows. Let $\cA_1:\cD\mapsto\cS_1$ be
$\eps_1$-DP, and let $\cA_2:\cD\times\cS_1\mapsto\cS_2$ be such that for every
fixed $y\in\cS_1$ the map $D\mapsto\cA_2(D,y)$ is $\eps_2$-DP. Define
$\cA(D)=(Y_1,Y_2)$ with $Y_1=\cA_1(D)$ and $Y_2=\cA_2(D,Y_1)$, where the
internal randomness of $\cA_1$ and $\cA_2$ is independent. Prove that $\cA$ is
$(\eps_1+\eps_2)$-DP. You may assume $\cS_1$ and $\cS_2$ are finite (or work
with densities).

\begin{answer}

\end{answer}

\item (5 pts) \textbf{The privacy region.} Consider the membership inference
game of Lecture 1 with $H_0$: the mechanism was run on $D$ and $H_1$: it was
run on the adjacent $D'$. An attacker is any (possibly randomized) test $T$
that maps the output to $\{0,1\}$; write $\FPR=\Pr(T=1\mid H_0)$ and
$\TPR=\Pr(T=1\mid H_1)$. Prove that if $\cA$ is $(\eps,\delta)$-DP then every
attacker satisfies
\begin{equation}\label{eq:privacy-region}
  \TPR\le e^{\eps}\,\FPR+\delta
  \qquad\text{and}\qquad
  1-\FPR\le e^{\eps}(1-\TPR)+\delta.
\end{equation}
\textit{Hint.} For a deterministic test, take $S=\{y:T(y)=1\}$ and its
complement; then handle randomized tests using part (b).

\begin{answer}

\end{answer}

\item (3 pts) \textbf{Reading the numbers.} A DP-trained model is released
with $(\eps,\delta)=(8,10^{-5})$, a common setting in practice. Using
\eqref{eq:privacy-region}, what is the largest $\TPR$ an attacker can have at
$\FPR=10^{-3}$? Repeat for $\eps=1$. Then explain in one or two sentences why
an attacker that always outputs 1 (so $\TPR=1$) does not contradict either
bound.

\begin{answer}

\end{answer}

\end{enumerate}

% ===============================================================
\problem{Problem 2: Understanding the basic mechanisms}{25 points}
% ===============================================================

Each part below examines a common misunderstanding about the Laplace,
Gaussian, or randomized response mechanism.

\begin{enumerate}

\item (6 pts) \textbf{Wrong norm, wrong guarantee.} A student implements the
Laplace mechanism for $f:\cD\mapsto\R^d$ but calibrates the noise to the
$\ell_2$ sensitivity: they release $f(D)+Z$ with
$Z_i\overset{\text{iid}}{\sim}\Lap(0,\Delta_2 f/\eps)$. Show that this
mechanism is $\sqrt{d}\,\eps$-DP, and give a function $f$ (with its adjacency)
for which this factor $\sqrt{d}$ cannot be improved, i.e.\ the mechanism is
not $\eps'$-DP for any $\eps'<\sqrt{d}\,\eps$.

\textit{Hint.} $\|x\|_1\le\sqrt{d}\,\|x\|_2$, with equality when all
coordinates have the same magnitude.

\begin{answer}

\end{answer}

\item (6 pts) \textbf{The Gaussian mechanism is never pure DP.} Take $d=1$,
$\Delta_2 f=\Delta>0$, and any $\sigma>0$. Show that the Gaussian mechanism
$f(D)+\cN(0,\sigma^2)$ is not $(\eps,0)$-DP for any finite $\eps$. Explain in
words why the Laplace distribution avoids this problem, and why the theorem
for the Gaussian mechanism therefore must have $\delta>0$ no matter how large
$\sigma$ is.

\begin{answer}

\end{answer}

\item (6 pts) \textbf{The strongest attacker against randomized response.}
Let $X\in\{0,1\}$ be one person's bit and $Y=M_{\mathrm{RR}}(X)$ the output of
randomized response with $p=e^{\eps}/(1+e^{\eps})$. Consider the hypotheses
$H_0$: $X=0$ and $H_1$: $X=1$.
\begin{enumerate}
  \item Compute the likelihood ratio
  $\Pr(Y=y\mid H_1)/\Pr(Y=y\mid H_0)$ for $y\in\{0,1\}$, and write down the
  Neyman--Pearson (likelihood ratio) test.

  \begin{answer}

  \end{answer}

  \item Compute the $(\FPR,\TPR)$ of the test ``output 1 if and only if
  $Y=1$'', and show that it satisfies the first inequality of
  \eqref{eq:privacy-region} with equality (with $\delta=0$). What does this
  say about the tightness of the DP bound for randomized response?

  \begin{answer}

  \end{answer}
\end{enumerate}

\item (7 pts) \textbf{Randomized response versus the Laplace mechanism.} We
want the fraction $\mu=\frac{1}{n}\sum_{i=1}^n X_i$ of $n$ bits
$X_i\in\{0,1\}$ with $\eps$-DP. Two options:
\begin{description}[leftmargin=4.2em, style=nextline]
  \item[(RR)] Each person applies $M_{\mathrm{RR}}$ to their own bit with
  parameter $\eps$; the analyst uses the unbiased estimator $\hat\mu$ from
  the lecture.
  \item[(Lap)] A trusted curator computes $\mu$ and releases
  $\mu+\Lap(0,\Delta_1\mu/\eps)$ under replacement adjacency.
\end{description}
Compute $\V[\hat\mu]$ exactly (not just the upper bound $\frac14$ on $p(1-p)$
used in the lecture) and $\V[\mu+Z]$ for (Lap). Show that the ratio of the two
variances is
\[
  \frac{V_{\mathrm{RR}}}{V_{\mathrm{Lap}}}
  =\frac{n\,\eps^2 e^{\eps}}{2\,(e^{\eps}-1)^2},
\]
which tends to $n/2$ as $\eps\to0$. Evaluate it for $n=1000$ and
$\eps\in\{0.1,1\}$. Then answer: the two options have the same $\eps$, so why
is (RR) so much noisier? Which of the two protects each person against a
stronger adversary, and in what sense? (One paragraph; we will return to this
in the lecture on local DP.)

\begin{answer}

\end{answer}

\end{enumerate}

% ===============================================================
\problem{Problem 3: Releasing a private feature mean (a DPSGD preview)}{25 points}
% ===============================================================

A hospital holds records of $n=1000$ patients. Each patient $i$ is described
by a feature vector $x_i\in\R^d$ with $d=100$, and the hospital wants to
release the mean $\bar{x}=\frac{1}{n}\sum_{i=1}^n x_i$ to an external research
group. (This is exactly the shape of a mini-batch gradient in DPSGD, which we
will study in Lecture 4: replace ``feature vector'' by ``per-example
gradient''.) Work under replacement adjacency, with $n$ public.

\begin{enumerate}

\item (5 pts) \textbf{Sensitivities.} Assume for now that every feature vector
satisfies $\|x_i\|_2\le1$. Compute $\Delta_2\bar{x}$. Then show
$\Delta_1\bar{x}\le\frac{2\sqrt{d}}{n}$ and exhibit a pair of vectors showing
this is tight.

\begin{answer}

\end{answer}

\item (6 pts) \textbf{Gaussian release.} The hospital targets
$(\eps,\delta)=(0.5,10^{-5})$. Using the Gaussian mechanism theorem from the
lecture, compute the required $\sigma$, the expected squared error
$\E\|Z\|_2^2$, and the root-mean-square error. Compare the latter to the
maximum possible size of the signal $\|\bar{x}\|_2$ and comment on whether the
release is useful. Why can you not use the same theorem to compute $\sigma$
for $\eps=2$, and what would you use instead?

\begin{answer}

\end{answer}

\item (5 pts) \textbf{Laplace release.} Repeat with the Laplace mechanism at
$\eps=0.5$ (pure DP). Show that the ratio of the expected squared errors,
Laplace over Gaussian, equals $d/\ln(1.25/\delta)$ in general, and evaluate it
here. In one sentence, state the trade-off the hospital is making when it
chooses Gaussian over Laplace.

\begin{answer}

\end{answer}

\item (5 pts) \textbf{Clipping.} In reality the features are not bounded. The
hospital's engineer proposes: ``compute $C=\max_i\|x_i\|_2$ from the data,
then release the Gaussian mechanism with $\Delta_2=2C/n$.'' Explain why the
resulting release is not $(\eps,\delta)$-DP. Then describe the standard fix
(clip every $x_i$ to $\ell_2$ norm at most a fixed $C$ before averaging, that
is, replace $x_i$ by $x_i\cdot\min(1,C/\|x_i\|_2)$), argue carefully why the
fix has bounded sensitivity, and describe the price paid in utility. Is the
clipped-then-noised mean an unbiased estimate of $\bar{x}$?

\begin{answer}

\end{answer}

\item (4 pts) \textbf{Repeated releases.} The hospital now wants to release an
updated mean every day for $T=100$ days from the same patient records (the
patients' features change from day to day, but the same $n$ people are in
every release). Using simple composition, what per-day budget $\eps_t$
achieves a total of $\eps=0.5$, and by what factor does $\sigma$ grow compared
to part (b)? What is the RMS error now? In one or two sentences, explain why
this calculation makes iterative algorithms such as DPSGD look hopeless, and
why the situation is not actually that bad (name the tools from later
lectures that improve on simple composition; you need not explain how they
work).

\begin{answer}

\end{answer}

\end{enumerate}

% ===============================================================
\problem{Problem 4: Record-level versus user-level DP on a social platform}{25 points}
% ===============================================================

A social platform has $N=10^6$ users. Each user $u$ has written some number
$k_u\ge0$ of posts over the past year. The platform's analyst wants to publish
statistics about posting activity with a DP guarantee. Two adjacency relations
are on the table:
\begin{itemize}
  \item \textbf{Record-level add/remove:} $D\simeq D'$ if $D'$ is obtained
  from $D$ by adding or removing one post.
  \item \textbf{User-level add/remove:} $D\simeq D'$ if $D'$ is obtained from
  $D$ by adding or removing one user together with all of their posts.
\end{itemize}

\begin{enumerate}

\item (5 pts) \textbf{Total number of posts.} Let $f(D)=\sum_u k_u$ be the
total number of posts. Compute $\Delta_1 f$ under record-level adjacency. Show
that under user-level adjacency $\Delta_1 f$ is unbounded. Then define the
contribution-bounded query $f_m(D)=\sum_u\min(k_u,m)$ and compute its
user-level sensitivity. What does $f_m$ measure when many users have
$k_u>m$?

\begin{answer}

\end{answer}

\item (6 pts) \textbf{Two routes to user-level pure DP.} The analyst wants
user-level $\eps_{\user}=5$ for every user, and knows that most users have at
most $m=50$ posts.
\begin{description}[leftmargin=5.5em, style=nextline]
  \item[Route A (convert)] Release $f(D)+\Lap(0,1/\eps_0)$ with record-level
  $\eps_0$, then invoke group privacy.
  \item[Route B (design)] Release $f_{50}(D)+\Lap(0,\Delta_1 f_{50}/\eps_{\user})$
  directly under user-level adjacency.
\end{description}
For Route A, what $\eps_0$ is needed so that a user with exactly 50 posts gets
$\eps_{\user}=5$? Compute the Laplace scale parameter under both routes and
show they are equal. Then explain why Route B still gives the strictly better
guarantee: what does Route A promise to a user with 200 posts, and what does
Route B promise to the same user?

\begin{answer}

\end{answer}

\item (6 pts) \textbf{The two routes under approximate DP.} Now the analyst
uses the Gaussian mechanism instead, with a record-level
$(\eps_0,\delta_0)=(0.1,10^{-8})$ in Route A, and user-level
$(\eps,\delta)=(5,10^{-8})$ on $f_{50}$ in Route B.
\begin{enumerate}
  \item Using the group-privacy theorem for approximate DP from the lecture,
  compute the user-level $(\eps_{\user},\delta_{\user})$ that Route A gives to
  a user with 50 posts.

  \begin{answer}

  \end{answer}

  \item Calibrate both routes with the noise formula
  $\sigma=\frac{\Delta_2 f}{\eps}\sqrt{2\ln(1.25/\delta)}$ at
  $\delta=10^{-8}$ (Route A: $\Delta_2 f=1$, $\eps=\eps_0$; Route B:
  $\Delta_2 f_{50}=50$, $\eps=5$). Show that the two values of $\sigma$
  coincide, so that for a user with exactly 50 posts the two routes release
  the same random variable. Route B certifies this mechanism as user-level
  $(5,10^{-8})$-DP for that user (the classical theorem does not apply at
  $\eps=5$, but you may take for granted, or verify numerically, that this
  $\sigma$ satisfies the analytic Gaussian mechanism condition at
  $(5,10^{-8})$). Explain why Route A's generic group-privacy conversion in
  (i) certifies a much worse $\delta$ for the very same mechanism, and state
  the lesson. Finally compare the $\delta_{\user}$ of (i) with the rule of
  thumb $\delta\le1/N$ from the lecture.

  \begin{answer}

  \end{answer}
\end{enumerate}

\item (5 pts) \textbf{A time series.} The analyst now wants to release the
vector $h\in\Z_{\ge0}^{365}$ of daily post counts, one coordinate per day of
the year, with pure $\eps$-DP.
\begin{enumerate}
  \item Under record-level adjacency, compute $\Delta_1 h$ and state the
  Laplace noise needed per coordinate. Compare with the naive plan ``treat the
  365 daily counts as 365 separate queries and split $\eps$ by simple
  composition'': by what factor is the naive noise larger, and why is the
  split unnecessary?

  \begin{answer}

  \end{answer}

  \item Under user-level adjacency, compute $\Delta_1 h$ for two ways of
  bounding contributions: (1) cap each user at $m$ posts in total for the
  year; (2) cap each user at $c$ posts per day. Which cap would you choose if
  most users write at most 3 posts per day but a few write 100, and why?
  Comment on what happens to the data of the heavy users under your choice.

  \begin{answer}

  \end{answer}
\end{enumerate}

\item (3 pts) \textbf{Writing the privacy statement.} The analyst's report
says only: ``All statistics are released with $\eps=1$.'' List what is
missing from this statement for a reader to know what is protected, and
explain why $\eps=1$ under record-level adjacency and $\eps=1$ under
user-level adjacency are not comparable promises, using your answer to
part (b).

\begin{answer}

\end{answer}

\end{enumerate}

\end{document}
